\documentclass[11pt]{article}

\usepackage[utf8]{inputenc}
\usepackage[T1]{fontenc}
\usepackage{lmodern}
\usepackage[margin=1in]{geometry}
\usepackage{amsmath, amssymb, amsthm, mathtools}
\usepackage{hyperref}
\hypersetup{colorlinks=true, linkcolor=blue, urlcolor=blue, citecolor=blue}

\newtheorem{theorem}{Theorem}
\newtheorem{lemma}[theorem]{Lemma}
\newtheorem{corollary}[theorem]{Corollary}
\theoremstyle{definition}
\newtheorem{definition}[theorem]{Definition}
\theoremstyle{remark}
\newtheorem{remark}[theorem]{Remark}

\title{ Permutation-aligned averaging: hidden-unit symmetry and the linear assignment reduction }
\author{Companion proof for: \href{https://github.com/pleyva2004/communication-efficient-learning-of-deep-networks}{pleyva2004/communication-efficient-learning-of-deep-networks} \\ \small Improvement slug: \texttt{permutation-aligned-averaging}}
\date{}

\begin{document}
\maketitle

\iffalse
NOTE: this file is a per-improvement proof, generated by the study-paper skill, Stage 7.
write-latex compliant: no fontspec, no XeLaTeX-only packages.
PROOF mode is selected when the improvement is mathematical, theoretical, or a
complexity-class statement (see Stage 7 heuristic). The proof must be complete:
no "obvious", no "left as exercise", no "by standard arguments".
This file loads NO natbib/bibliography; prior work is cited as inline plaintext only.
\fi

\begin{abstract}
The math deep dive (\S7b) attributes Figure~1's loss barrier between independently-initialized
networks to the permutation symmetry of a hidden layer, and improvement T.1 of
\texttt{05-improvements.tex} proposes to remove that barrier before averaging by re-aligning each
client model into a common basin. This note makes both halves rigorous. We prove (i) that for a
one-hidden-layer network $h(x)=W_2\,\sigma(W_1 x + b_1)$ with $\sigma$ applied elementwise, the
reparametrization $W_1\mapsto P W_1$, $b_1\mapsto P b_1$, $W_2\mapsto W_2 P^{\top}$ by \emph{any}
permutation matrix $P$ leaves the function $h$ \emph{exactly} unchanged, hence leaves any loss
unchanged; and (ii) that minimizing the Frobenius distance $\lVert W_1 - P W_1'\rVert_F^2$ over
permutation matrices is equivalent to the linear assignment problem of maximizing
$\sum_i \langle (W_1)_{i,:},\,(W_1')_{P(i),:}\rangle$, because the two squared-norm terms are
permutation-invariant and drop out --- a problem solvable exactly in $O(H^3)$ time by the Hungarian
algorithm. Part (i) follows from $P^{\top}P=I$ together with the fact that an elementwise
nonlinearity commutes with a coordinate permutation; part (ii) is a direct expansion of the
Frobenius inner product. Together they justify ``align then average'' (weight-matching, in the sense
of Git Re-Basin, Ainsworth et al. 2023) as an exact symmetry reduction with a tractable matching
step.
\end{abstract}

\section*{Setup and notation}

We work with a single-hidden-layer feed-forward network with $H$ hidden units. Fix an input
dimension $d_{\mathrm{in}}$, a hidden width $H$, and an output dimension $d_{\mathrm{out}}$. The
network is parametrized by a first-layer weight matrix $W_1\in\mathbb{R}^{H\times d_{\mathrm{in}}}$,
a hidden bias $b_1\in\mathbb{R}^{H}$, and a second-layer weight matrix
$W_2\in\mathbb{R}^{d_{\mathrm{out}}\times H}$. A fixed scalar activation
$\phi:\mathbb{R}\to\mathbb{R}$ (e.g.\ ReLU $\phi(t)=\max(t,0)$, or $\tanh$) is applied
\emph{elementwise}: for a vector $z=(z_1,\dots,z_H)^{\top}\in\mathbb{R}^{H}$ we write
$\sigma(z)=(\phi(z_1),\dots,\phi(z_H))^{\top}$, so $\sigma:\mathbb{R}^H\to\mathbb{R}^H$ acts
coordinatewise with the \emph{same} scalar map $\phi$ in every coordinate. The forward map is
\[
h(x)\;=\;W_2\,\sigma\!\bigl(W_1 x + b_1\bigr),\qquad x\in\mathbb{R}^{d_{\mathrm{in}}}.
\]
The deep dive's notation in \S7b uses $W_2\sigma(W_1x+b_1)$ for the 2NN with $200$ hidden units per
layer; here $H$ plays the role of that hidden width.

A \emph{permutation matrix} $P\in\mathbb{R}^{H\times H}$ is the matrix of a bijection
$\pi:\{1,\dots,H\}\to\{1,\dots,H\}$, defined by $P_{ij}=1$ if $j=\pi(i)$ and $P_{ij}=0$ otherwise;
equivalently $P_{ij}=\mathbf 1\{j=\pi(i)\}$. We write $P(\cdot)$ for the induced reparametrization of
the parameter triple and $\Pi_H$ for the set of all $H\times H$ permutation matrices, a finite group
of order $H!$ under matrix multiplication. We use the Frobenius inner product
$\langle A,B\rangle_F=\sum_{ij}A_{ij}B_{ij}=\operatorname{tr}(A^{\top}B)$ and the induced norm
$\lVert A\rVert_F^2=\langle A,A\rangle_F$. Rows are written $(W)_{i,:}$ and columns $(W)_{:,j}$. The
linear assignment problem (LAP) on a cost/benefit matrix $C\in\mathbb{R}^{H\times H}$ asks for the
bijection $\pi$ optimizing $\sum_i C_{i,\pi(i)}$.

\begin{definition}[Working notation]
For a permutation $\pi$ with matrix $P$ ($P_{ij}=\mathbf 1\{j=\pi(i)\}$):
\begin{itemize}
  \item \textbf{Action on a vector.} For $z\in\mathbb{R}^H$, the $i$-th coordinate of $Pz$ is
  $(Pz)_i=\sum_{j}P_{ij}z_j=z_{\pi(i)}$. Thus $P$ \emph{gathers}: it places old coordinate
  $\pi(i)$ into new slot $i$.
  \item \textbf{Orthogonality.} Each row and each column of $P$ has exactly one $1$, so the columns
  of $P$ are the standard basis vectors in some order; hence $P^{\top}P=PP^{\top}=I_H$, i.e.\
  $P^{-1}=P^{\top}$ and $P$ is orthogonal.
  \item \textbf{Elementwise nonlinearity.} $\sigma:\mathbb{R}^H\to\mathbb{R}^H$,
  $\sigma(z)_i=\phi(z_i)$, with the \emph{same} scalar $\phi$ in every coordinate.
  \item \textbf{Reparametrization.} $P$ acts on the parameter triple by
  $\;W_1\mapsto P W_1,\quad b_1\mapsto P b_1,\quad W_2\mapsto W_2 P^{\top}.$
  \item \textbf{Loss.} For any data distribution $\mathcal D$ over $(x,y)$ and any per-example loss
  $\ell$, the population (or empirical) risk of a parameter set $\theta=(W_1,b_1,W_2)$ is
  $L(\theta)=\mathbb{E}_{(x,y)\sim\mathcal D}\,\ell\bigl(h_\theta(x),y\bigr)$, where $h_\theta$ is the
  forward map above.
\end{itemize}
\end{definition}

\section*{Theorem}

\begin{theorem}[Hidden-unit permutation symmetry and the assignment reduction]
\label{thm:main}
Let $h(x)=W_2\,\sigma(W_1 x + b_1)$ be the one-hidden-layer network of the Setup, with $\sigma$
applied elementwise via a fixed scalar $\phi$. Then:

\smallskip
\noindent\textbf{(i) Exact functional symmetry.} For every permutation matrix $P\in\Pi_H$, the
reparametrized network
\[
\widetilde W_1=P W_1,\qquad \widetilde b_1 = P b_1,\qquad \widetilde W_2 = W_2 P^{\top}
\]
computes the identical function: $\widetilde W_2\,\sigma(\widetilde W_1 x + \widetilde b_1)=h(x)$ for
\emph{all} $x\in\mathbb{R}^{d_{\mathrm{in}}}$. Consequently, for any data distribution and any loss
$\ell$, the risk is unchanged: $L\bigl(P W_1, P b_1, W_2 P^{\top}\bigr)=L\bigl(W_1,b_1,W_2\bigr)$.

\smallskip
\noindent\textbf{(ii) Frobenius alignment $\equiv$ linear assignment.} Let
$W_1,W_1'\in\mathbb{R}^{H\times d_{\mathrm{in}}}$ be two first-layer weight matrices. Then for every
$P\in\Pi_H$ with associated permutation $\pi$,
\[
\bigl\lVert W_1 - P\,W_1'\bigr\rVert_F^2
\;=\;\lVert W_1\rVert_F^2 + \lVert W_1'\rVert_F^2
\;-\;2\sum_{i=1}^{H}\bigl\langle (W_1)_{i,:},\,(W_1')_{\pi(i),:}\bigr\rangle ,
\]
where the first two terms are constants independent of $P$. Hence
\[
\arg\min_{P\in\Pi_H}\bigl\lVert W_1 - P\,W_1'\bigr\rVert_F^2
\;=\;\arg\max_{P\in\Pi_H}\;\sum_{i=1}^{H}\bigl\langle (W_1)_{i,:},\,(W_1')_{\pi(i),:}\bigr\rangle .
\]
The right-hand side is a linear assignment problem on the $H\times H$ benefit matrix
$C_{i,j}=\langle (W_1)_{i,:},\,(W_1')_{j,:}\rangle$, and it can be solved \emph{exactly} in
$O(H^3)$ time (the Hungarian algorithm of Kuhn 1955 and Munkres 1957).
\end{theorem}

\section*{Proof}

\begin{proof}
We prove the two parts in turn. Throughout, $P\in\Pi_H$ is the matrix of a bijection $\pi$ with
$P_{ij}=\mathbf 1\{j=\pi(i)\}$.

\medskip
\noindent\textbf{Part (i): exact functional symmetry.}

\emph{Step 1 --- $P$ is orthogonal, $P^{\top}P=I_H$.}
Fix indices $i,k\in\{1,\dots,H\}$ and compute the $(i,k)$ entry of $P^{\top}P$. By definition of
matrix multiplication and of $P$,
\[
\bigl(P^{\top}P\bigr)_{ik}
=\sum_{m=1}^{H} (P^{\top})_{im}\,P_{mk}
=\sum_{m=1}^{H} P_{mi}\,P_{mk}
=\sum_{m=1}^{H} \mathbf 1\{i=\pi(m)\}\,\mathbf 1\{k=\pi(m)\}.
\]
Each summand is $1$ exactly when $\pi(m)=i$ and $\pi(m)=k$ simultaneously, which requires $i=k$.
Because $\pi$ is a bijection, for each $i$ there is exactly one index $m=\pi^{-1}(i)$ with
$\pi(m)=i$. Hence the sum equals $1$ if $i=k$ and $0$ otherwise, i.e.\
$\bigl(P^{\top}P\bigr)_{ik}=\delta_{ik}$, so $P^{\top}P=I_H$. (The same computation with the roles of
rows and columns swapped gives $PP^{\top}=I_H$; we only need $P^{\top}P=I_H$ below.) In particular
$P^{\top}=P^{-1}$.

\emph{Step 2 --- an elementwise nonlinearity commutes with a coordinate permutation.}
We claim that for every $z\in\mathbb{R}^H$,
\[
\sigma\bigl(P z\bigr)=P\,\sigma\bigl(z\bigr).
\]
Both sides are vectors in $\mathbb{R}^H$; we compare their $i$-th coordinates. Recall from the
working notation that $(Pz)_i=\sum_j P_{ij}z_j = z_{\pi(i)}$, since the only nonzero entry in row $i$
of $P$ is $P_{i,\pi(i)}=1$. Applying $\sigma$ coordinatewise,
\[
\bigl(\sigma(Pz)\bigr)_i=\phi\bigl((Pz)_i\bigr)=\phi\bigl(z_{\pi(i)}\bigr).
\]
On the other side, $\sigma(z)$ is the vector with $j$-th coordinate $\phi(z_j)$, and applying $P$ to
it gathers coordinate $\pi(i)$ into slot $i$:
\[
\bigl(P\,\sigma(z)\bigr)_i=\sum_{j}P_{ij}\,\sigma(z)_j=\sigma(z)_{\pi(i)}=\phi\bigl(z_{\pi(i)}\bigr).
\]
The two expressions agree for every $i$, so $\sigma(Pz)=P\sigma(z)$. The only property of $\sigma$
used is that it acts coordinatewise with one and the same scalar $\phi$ in every coordinate; this is
exactly what ``$\sigma$ applied elementwise'' means, and it is what makes the permutation slide
through the nonlinearity. (If different coordinates used different scalar functions, this step would
fail.)

\emph{Step 3 --- assemble the forward pass.}
Let $\widetilde W_1=PW_1$, $\widetilde b_1=Pb_1$, $\widetilde W_2=W_2P^{\top}$, and fix an arbitrary
input $x\in\mathbb{R}^{d_{\mathrm{in}}}$. The reparametrized pre-activation is
\[
\widetilde W_1 x + \widetilde b_1 = (PW_1)x + (Pb_1) = P(W_1 x) + P b_1 = P\bigl(W_1 x + b_1\bigr),
\]
using only distributivity of matrix multiplication and associativity $(PW_1)x=P(W_1x)$. Write
$z:=W_1 x + b_1\in\mathbb{R}^H$ for the original pre-activation, so the reparametrized
pre-activation is $Pz$. Pushing through $\sigma$ and then the second layer:
\[
\widetilde W_2\,\sigma\bigl(\widetilde W_1 x + \widetilde b_1\bigr)
= (W_2 P^{\top})\,\sigma\bigl(P z\bigr)
\overset{\text{Step 2}}{=} (W_2 P^{\top})\,\bigl(P\,\sigma(z)\bigr)
= W_2\,\bigl(P^{\top}P\bigr)\,\sigma(z)
\overset{\text{Step 1}}{=} W_2\,\sigma(z),
\]
where the middle equality used associativity to regroup $(W_2P^{\top})(P\sigma(z))=W_2(P^{\top}P)\sigma(z)$,
and the last used $P^{\top}P=I_H$ from Step~1. But $W_2\,\sigma(z)=W_2\,\sigma(W_1x+b_1)=h(x)$ by
definition. Since $x$ was arbitrary,
\[
\widetilde W_2\,\sigma\bigl(\widetilde W_1 x + \widetilde b_1\bigr)=h(x)\qquad\text{for all }x,
\]
i.e.\ the two parameter triples realize the \emph{same} function $\mathbb{R}^{d_{\mathrm{in}}}\to
\mathbb{R}^{d_{\mathrm{out}}}$.

\emph{Step 4 --- loss invariance.}
Loss depends on the parameters only through the function they compute, because each example's loss
$\ell(h_\theta(x),y)$ is a function of the network \emph{output} $h_\theta(x)$. Let
$\theta=(W_1,b_1,W_2)$ and $\widetilde\theta=(PW_1,Pb_1,W_2P^{\top})$. By Step~3,
$h_{\widetilde\theta}(x)=h_\theta(x)$ for every $x$, so $\ell(h_{\widetilde\theta}(x),y)=
\ell(h_\theta(x),y)$ for every $(x,y)$. Taking expectation over $(x,y)\sim\mathcal D$ (or the finite
empirical average over a fixed sample, which is the special case where $\mathcal D$ is the empirical
distribution) and using monotonicity/linearity of expectation on a pointwise identity,
\[
L(\widetilde\theta)=\mathbb{E}_{(x,y)}\,\ell\bigl(h_{\widetilde\theta}(x),y\bigr)
=\mathbb{E}_{(x,y)}\,\ell\bigl(h_{\theta}(x),y\bigr)=L(\theta).
\]
Hence $L(PW_1,Pb_1,W_2P^{\top})=L(W_1,b_1,W_2)$, proving (i). As a structural corollary, the map
$\theta\mapsto P(\theta)$ permutes the parameter space within a single level set of $L$; ranging over
all $P\in\Pi_H$ produces an orbit of (generically) $H!$ distinct parameter points that all attain the
identical loss, matching the deep dive's $(200!)^2$ orbit count for the two-hidden-layer 2NN
(one factor of $H!$ per hidden layer).

\medskip
\noindent\textbf{Part (ii): Frobenius alignment reduces to linear assignment.}

\emph{Step 1 --- expand the Frobenius norm.}
For any matrices $A,B$ of the same shape, $\lVert A-B\rVert_F^2=\langle A-B,A-B\rangle_F$. Expanding
the inner product by bilinearity,
\[
\lVert A-B\rVert_F^2
=\langle A,A\rangle_F-\langle A,B\rangle_F-\langle B,A\rangle_F+\langle B,B\rangle_F
=\lVert A\rVert_F^2 - 2\langle A,B\rangle_F + \lVert B\rVert_F^2,
\]
where we used symmetry $\langle A,B\rangle_F=\langle B,A\rangle_F$ of the (real) Frobenius inner
product. Apply this with $A=W_1$ and $B=PW_1'$:
\[
\bigl\lVert W_1 - P W_1'\bigr\rVert_F^2
=\lVert W_1\rVert_F^2 - 2\bigl\langle W_1,\,P W_1'\bigr\rangle_F + \bigl\lVert P W_1'\bigr\rVert_F^2.
\tag{$\ast$}
\]

\emph{Step 2 --- the third term is permutation-invariant.}
Using $\lVert M\rVert_F^2=\operatorname{tr}(M^{\top}M)$ and the cyclic property of the trace, with
$P^{\top}P=I_H$ from Part~(i) Step~1,
\[
\bigl\lVert P W_1'\bigr\rVert_F^2
=\operatorname{tr}\!\bigl((PW_1')^{\top}(PW_1')\bigr)
=\operatorname{tr}\!\bigl(W_1'^{\top}\,P^{\top}P\,W_1'\bigr)
=\operatorname{tr}\!\bigl(W_1'^{\top}W_1'\bigr)
=\lVert W_1'\rVert_F^2 .
\]
(Geometrically: $P$ is orthogonal, so left-multiplication by $P$ is an isometry and preserves the
Frobenius norm; concretely, the rows of $PW_1'$ are exactly the rows of $W_1'$ in the order
prescribed by $\pi$, and reordering rows does not change the sum of squared entries.) The first term
$\lVert W_1\rVert_F^2$ does not involve $P$ at all. Therefore in ($\ast$) both squared-norm terms are
constants independent of $P$, and
\[
\bigl\lVert W_1 - P W_1'\bigr\rVert_F^2
=\underbrace{\lVert W_1\rVert_F^2 + \lVert W_1'\rVert_F^2}_{\text{constant in }P}
\;-\;2\,\bigl\langle W_1,\,P W_1'\bigr\rangle_F .
\tag{$\ast\ast$}
\]

\emph{Step 3 --- rewrite the cross term as a sum of row inner products.}
Expand the Frobenius inner product entrywise and substitute $(P W_1')_{ik}=\sum_m P_{im}(W_1')_{mk}
=(W_1')_{\pi(i),k}$ (only $m=\pi(i)$ contributes):
\[
\bigl\langle W_1,\,P W_1'\bigr\rangle_F
=\sum_{i=1}^{H}\sum_{k=1}^{d_{\mathrm{in}}} (W_1)_{ik}\,(PW_1')_{ik}
=\sum_{i=1}^{H}\sum_{k=1}^{d_{\mathrm{in}}} (W_1)_{ik}\,(W_1')_{\pi(i),k}
=\sum_{i=1}^{H}\bigl\langle (W_1)_{i,:},\,(W_1')_{\pi(i),:}\bigr\rangle ,
\]
where the inner sum over $k$ is precisely the Euclidean inner product of row $i$ of $W_1$ with row
$\pi(i)$ of $W_1'$. Substituting into ($\ast\ast$) gives the claimed identity
\[
\bigl\lVert W_1 - P W_1'\bigr\rVert_F^2
=\lVert W_1\rVert_F^2 + \lVert W_1'\rVert_F^2
-2\sum_{i=1}^{H}\bigl\langle (W_1)_{i,:},\,(W_1')_{\pi(i),:}\bigr\rangle .
\]

\emph{Step 4 --- minimization becomes maximization of the assignment objective.}
For each fixed pair $(W_1,W_1')$ the two squared-norm terms are a single additive constant
$\kappa:=\lVert W_1\rVert_F^2 + \lVert W_1'\rVert_F^2$ that does not depend on the choice of
$P\in\Pi_H$. A function of the form $P\mapsto \kappa - 2\,S(P)$, with $\kappa$ constant and the
coefficient $-2<0$, is minimized exactly where $S(P):=\sum_i\langle (W_1)_{i,:},(W_1')_{\pi(i),:}
\rangle$ is maximized: for any $P,Q\in\Pi_H$,
\[
\bigl\lVert W_1 - P W_1'\bigr\rVert_F^2 \le \bigl\lVert W_1 - Q W_1'\bigr\rVert_F^2
\;\Longleftrightarrow\;
\kappa - 2S(P)\le \kappa - 2S(Q)
\;\Longleftrightarrow\;
S(P)\ge S(Q).
\]
Hence the set of minimizers of the left objective over the finite set $\Pi_H$ equals the set of
maximizers of $S$:
\[
\arg\min_{P\in\Pi_H}\bigl\lVert W_1 - P W_1'\bigr\rVert_F^2
=\arg\max_{P\in\Pi_H}\;\sum_{i=1}^{H}\bigl\langle (W_1)_{i,:},\,(W_1')_{\pi(i),:}\bigr\rangle .
\]
(Both $\arg\min$ and $\arg\max$ are over the finite set $\Pi_H$, so they are attained and nonempty;
the equivalence is of the full minimizer/maximizer sets, not merely of a single optimizer.)

\emph{Step 5 --- this is a linear assignment problem, solvable in $O(H^3)$.}
Define the benefit matrix $C\in\mathbb{R}^{H\times H}$ by
$C_{ij}=\langle (W_1)_{i,:},\,(W_1')_{j,:}\rangle$, i.e.\ $C=W_1\,W_1'^{\top}$ (entry $(i,j)$ is the
inner product of row $i$ of $W_1$ with row $j$ of $W_1'$). A permutation $\pi$ is a perfect matching
of the bipartite graph with left vertices $\{1,\dots,H\}$ (rows of $W_1$), right vertices
$\{1,\dots,H\}$ (rows of $W_1'$), and edge weight $C_{i,j}$. The objective $S(P)=\sum_{i}C_{i,\pi(i)}$
is exactly the total weight of the matching $\pi$, so maximizing $S$ over $\Pi_H$ is the
\emph{maximum-weight perfect bipartite matching} (equivalently the linear assignment) problem on $C$.
Maximization is converted to the standard minimization form by minimizing
$C'_{ij}:=\bigl(\max_{a,b}C_{ab}\bigr)-C_{ij}\ge 0$, an affine, $\pi$-independent transformation of
the per-edge weights that adds the constant $H\cdot\max_{a,b}C_{ab}$ to every assignment's total and
hence preserves the optimizing $\pi$. The Hungarian algorithm (Kuhn 1955; Munkres 1957) solves this
to global optimality --- it returns a provably cost-minimizing assignment, not a heuristic --- in
$O(H^3)$ arithmetic operations on the $H\times H$ matrix $C'$. Therefore the Frobenius-optimal
alignment $P^{\star}$ exists, is exactly computable, and is found in cubic time in the hidden width.
This proves (ii). \qedhere
\end{proof}

\begin{remark}[Aligning the whole network, not just $W_1$]
By Part~(i), applying $P^{\star}$ to the \emph{whole} second model --- $W_1'\mapsto P^{\star}W_1'$,
$b_1'\mapsto P^{\star}b_1'$, $W_2'\mapsto W_2'(P^{\star})^{\top}$ --- leaves the second model's
function (and loss) unchanged while making its hidden units maximally row-aligned with the first
model's. The server may then average the aligned parameters, e.g.\
$\tfrac12\theta + \tfrac12 P^{\star}(\theta')$, blending hidden units that play matching roles rather
than functionally unrelated ones. The matching above uses only $W_1$; one may instead build $C$ from
$[\,W_1\,|\,b_1\,]$ or incorporate $W_2$ columns, which changes the cost matrix but not the
reduction --- it remains a single $O(H^3)$ assignment.
\end{remark}

\section*{Discussion}

\paragraph{What this gives us.}
Part (i) is an \emph{exact} statement, holding for every $x$ and every permutation, with no
approximation, no smoothness assumption on $\phi$ beyond being a fixed scalar map, and no convexity:
hidden-unit relabeling is a genuine symmetry of the loss surface, so every minimizer sits on an
orbit of (generically) $H!$ equal-loss copies. This is the precise mechanism the math deep dive
invokes in \S7b to explain Figure~1's loss barrier between independently-initialized models: two
solutions found from different seeds generically land in \emph{different} elements of the same orbit,
so coordinate-wise averaging blends unrelated hidden units and produces a degraded interpolant. Part
(ii) turns the proposed fix into something tractable and exactly solvable: ``find the best
relabeling'' is not a search over $H!$ permutations but a polynomial-time linear assignment, because
the permutation-invariant squared-norm terms drop out of the Frobenius objective, leaving a pure
inner-product matching. Concretely this licenses the T.1 ``align then average'' procedure
(weight-matching, in the sense of Git Re-Basin, Ainsworth et al.\ 2023): solve one $O(H^3)$
Hungarian assignment on $C=W_1 W_1'^{\top}$, apply the resulting symmetry (which by Part~(i) costs
nothing in function value), then average. It is what relaxes FedAvg's reliance on a \emph{shared}
per-round initialization (McMahan et al.\ 2017): the shared init was the cheap way to keep all
clients in one orbit element, and explicit alignment provides an alternative that does not require it,
which in principle permits larger local epochs $E$ and fewer synchronization rounds in the
drift-prone non-IID regime that motivates client-drift corrections such as SCAFFOLD
(Karimireddy et al.\ 2020) and FedProx (Li et al.\ 2020).

\paragraph{What this does not give us.}
The theorem is purely about the \emph{symmetry} and the \emph{matching}; it says nothing about
whether averaging two aligned models is actually good. Part (i) guarantees only that alignment is
loss-preserving for the model being aligned, not that the post-alignment \emph{average} has low loss
--- that is an empirical, landscape-dependent claim (the prototype's barrier $+0.024\to-0.023$ in
T.1 is evidence, not a consequence of this proof). The matching in (ii) is optimal for the chosen
\emph{surrogate} objective $\lVert W_1 - PW_1'\rVert_F^2$ (raw first-layer weight distance); minimizing
weight distance is a proxy for, and not the same as, minimizing the post-average loss barrier, so the
``optimal'' $P^{\star}$ is optimal for the surrogate only. Nothing here addresses multi-client
aggregation: for $K>2$ clients, jointly aligning all of them to a common reference so that pairwise
matchings are mutually consistent is a separate combinatorial problem that the single $O(H^3)$
solve does not settle. Finally, the result is exact for the one-hidden-layer architecture stated; it
does not by itself extend the function-preservation guarantee to convolutions, normalization layers,
attention, or residual connections, each of which has its own (sometimes more constrained) symmetry
group.

\paragraph{Where it would break.}
Part (i)'s Step~2 is the load-bearing hinge and it requires the nonlinearity to be \emph{truly
elementwise with the same scalar in every coordinate}: a per-channel parametrized activation (e.g.\
PReLU with distinct learned slopes per unit), a softmax or any cross-coordinate coupling, a
LayerNorm/BatchNorm whose statistics mix coordinates, or a bias folded in non-permutably would all
break $\sigma(Pz)=P\sigma(z)$, and the symmetry would no longer be exact (it would have to be carried
by the matching layer's statistics too). The clean second-layer cancellation also depends on the
convention $W_2\mapsto W_2 P^{\top}$ with the \emph{same} $P$ that acts on the first layer; mismatched
or untied permutations would not produce $P^{\top}P=I$. For \emph{deeper} networks, weight-matching
is per-layer and greedy: aligning layer $\ell$ fixes a permutation that must be propagated to layer
$\ell+1$, and choosing each layer's permutation by its own local Frobenius objective is not jointly
optimal across layers --- this is why activation-matching (aligning units by their responses on real
data rather than by raw weight inner products) is generally needed for multilayer nets, and even then
the joint problem is only solved approximately by coordinate descent over per-layer assignments
(Ainsworth et al.\ 2023). Finally, the procedure is at odds with secure aggregation
(Bonawitz et al.\ 2017): computing $C=W_1 W_1'^{\top}$ and the matching requires the server to see
individual client weight matrices, whereas secure aggregation is designed precisely so the server
learns only the \emph{sum} of updates and never any single client's parameters; making alignment
compatible with that privacy model (e.g.\ matching to a public reference, or matching under
encryption) is a genuinely open, non-trivial obstacle, not a detail this proof resolves.

\paragraph{Empirical companion.}
This proof is the validation file for the improvement \texttt{permutation-aligned-averaging} in
\texttt{05-improvements.tex} (proposal T.1). If a measurement-mode prototype also exists at
\texttt{improvements/permutation-aligned-averaging.py}, it serves as a sanity check, not as evidence
--- the theorem above is the load-bearing claim.

\end{document}
